% Task 2 contribution. Stable IDs NX-L01--NX-L07.
% Whole-proof readings and their exact scope are recorded in CHECK_EVENTS.tsv.

\section{Bad fibres of $\Lambda$-schemes}

\subsection{The separatedness qualification is substantive}

A $\Lambda$-structure on a flat scheme over $\mathbb Z$ means commuting
endomorphisms $\psi_q$ lifting Frobenius at every prime $q$; the definition
itself does not require separatedness (Borger, introduction, p.~2, and
\S1.1, p.~6).  Borger explicitly adds separatedness at the beginning of
\S6 before Theorem~6.1 and Corollary~6.2.  The conjecture in the present
prompt says ``flat and of finite type'' without adding separatedness.
That literal statement has a counterexample.

\begin{theorem}[NX-L01: arbitrary extra bad factors without separatedness]
\label{NX-L01}
Fix a prime $p$, let $Y=\mathbb A^d_{\mathbb Z}$, and let
$Z\subseteq Y_{\mathbb F_p}$ be any closed subscheme.  Glue two copies
$Y_1,Y_2$ along the identity of $U=Y\setminus |Z|$, and call the resulting
scheme $X$.  Then $X$ is flat and of finite type over $\mathbb Z$ and
admits a $\Lambda$-structure.  Its generic fibre is $Y_{\mathbb Q}$ and
\[
 \zeta_X(s)=\zeta(s-d)\,Z(Z,p^{-s}).
\]
Thus its $p$-factor in the correction of ZETA\_TWO~ZT04 is exactly
$Z(Z,p^{-s})$.  In particular the proposed cyclotomic integral-Tate
factor assertion is false as stated.
\end{theorem}

\begin{proof}
Write $\varphi_q:Y\to Y$ for coordinatewise $q$th powers.  The fold map
$r:X\to Y$ is the identity on each chart, and write $i_1:Y\to X$ for the
first chart.  Because Frobenius on $Y_{\mathbb F_p}$ is the identity on
its topological space, $\varphi_p^{-1}(U)=U$.  Consequently the two
maps $\varphi_p:Y_a\to Y_a$ glue to an endomorphism $\psi_p$ of $X$.
For $q\ne p$, set
\[
 \psi_q=i_1\circ\varphi_q\circ r.
\]
The identities $r i_1=1$, $r\psi_p=\varphi_p r$ and
$\psi_p i_1=i_1\varphi_p$, together with
$\varphi_q\varphi_{q'}=\varphi_{q'}\varphi_q$, prove pairwise
commutativity for all primes.  At $q\ne p$ the overlap is the whole
fibre, so $r:X_{\mathbb F_q}\simeq Y_{\mathbb F_q}$ and $\psi_q$
restricts to its absolute Frobenius.  At $p$ this assertion holds on
both affine charts by construction.  These commuting Frobenius lifts
supply the $\Lambda$-structure by the flat criterion in Borger~\S1.1.
Flatness and local finite type are local on the source; the two affine
charts are flat and of finite type, and their finite union is
quasi-compact.  Thus $X$ satisfies the hypotheses in the assertion.

For each $r\geq1$, the two copies have precisely the points of
$Y(\mathbb F_{p^r})\setminus Z(\mathbb F_{p^r})$ identified.  Hence
\[
 \#X(\mathbb F_{p^r})=p^{dr}+\#Z(\mathbb F_{p^r}).
\]
All other fibres are those of $Y$.  Taking the exponential of
$\sum_{r\geq1}\#X(\mathbb F_{p^r})t^r/r$ proves the local identity,
then the global Euler product and the stated ZT04 correction.

For a completely explicit choice, take $p=5$, $d=2$, and
$Z:\ y^2=x^3-x$ in $\mathbb A^2_{\mathbb F_5}$.  Its smooth projective
completion $E$ (discriminant $64\ne0$ in $\mathbb F_5$) has one
point at infinity and eight points over $\mathbb F_5$: the five affine $x$-values give respectively
$1,1,2,2,1$ choices for $y$.  Silverman, \emph{The Arithmetic of Elliptic Curves}, second edition,
Theorem~V.2.4 (p.~143), gives
\[
 Z(E,t)=\frac{1+2t+5t^2}{(1-t)(1-5t)},\qquad
 Z(Z,t)=\frac{1+2t+5t^2}{1-5t}.
\]
The two numerator inverse roots are $-1+2i,-1-2i$, of absolute value
$\sqrt5$, whereas every inverse root of a product of factors
$(1-\xi5^jt)^{\pm1}$ with $j\in\mathbb Z$ has absolute value an
integral power of $5$.  The numerator has no root in common with
$1-5t$, so cancellation cannot repair the discrepancy.  This proves
refutation.  The chosen $X$ is not separated: the two sections of $X$
obtained by lifting the point $(0,0)$ agree over $\mathbb Z[1/5]$ and
disagree at $5$.  In particular this example says nothing against a
version that expressly assumes separatedness.
\end{proof}

\subsection{A local comparison for a proper Frobenius lift}

For the proper question, a comparison must go beyond proper base
change: proper base change alone does not identify the generic and
special fibres of a singular family.  The following construction uses
the lift at the bad prime.  Its finite-coefficient form is recorded
first because filtered colimits and passage to $\mathbb Q_\ell$ must
not be interchanged without an argument.

\begin{theorem}[NX-L02: the proper Frobenius tower]
\label{NX-L02}
Let $\mathcal X$ be a proper flat scheme of finite type over
$\mathbb Z_p$, and let $\varphi:\mathcal X\to\mathcal X$ be a
$\mathbb Z_p$-endomorphism whose special fibre is absolute Frobenius.
For every prime $\ell\ne p$ and every $a\geq1$, there are natural,
$G_{\mathbb Q_p}$-equivariant isomorphisms
\[
 H^i(\mathcal X_{\overline{\mathbb F}_p},\mathbb Z/\ell^a)
 \simeq
 \mathop{\rm colim}_{\varphi^*}
 H^i(\mathcal X_{\overline{\mathbb Q}_p},\mathbb Z/\ell^a).
 \tag{L2.1}
\]
Here $G_{\mathbb Q_p}$ acts on the left through its unramified
quotient.  In particular
$H^i(\mathcal X_{\overline{\mathbb F}_p},\mathbb Q_\ell)$ is
isomorphic to a $G_{\mathbb Q_p}$-stable direct summand of
$H^i(\mathcal X_{\overline{\mathbb Q}_p},\mathbb Q_\ell)$.
\end{theorem}

\begin{proof}
Put $C=\widehat{\overline{\mathbb Q}_p}$ and let $C^\flat$ be its tilt.
We use the following precise analytic results: Huber (1994),
Remark~4.6(iv)(d), identifies the generic fibre of the formal completion
of a proper scheme with the analytification of its algebraic generic
fibre; Scholze, \emph{Perfectoid spaces}, Theorem~5.2 gives tilting of
perfectoid algebras, Theorem~7.12 gives the equivalence of their
\'etale sites, Definition~7.14 and Corollary~7.18 give cohomological
continuity for perfectoid tilde-limits, and Theorem~9.7 (Huber's
Theorem~3.8.1) gives the algebraic--adic comparison with the finite
coefficients used here.  These are cited, not proved again below.
The construction required to apply them in this setting is as follows.

Complete $\mathcal X$ along its special fibre.  The topological
space of the resulting formal scheme $\mathfrak X$ is
$|\mathcal X_{\mathbb F_p}|$, on which $\varphi$ is the identity.
Every formal affine open $\mathfrak U=\operatorname{Spf}A$ is
therefore invariant under inverse image by $\varphi$, and the map on
it is a continuous endomorphism of $A$ lifting Frobenius modulo $p$.
Set
\[
 B=(A/pA)^{\mathrm{perf}},\qquad
 A_\infty=\widehat{\mathop{\rm colim}_{\varphi} A}^{\,(p)}.
 \tag{L2.2}
\]
The ring $A$ is $p$-torsion-free, by flatness.  Filtered colimits of
$p$-torsion-free $\mathbb Z_p$-modules are $p$-torsion-free.  For any
$p$-torsion-free ring $D$, its $p$-adic completion is again
$p$-torsion-free and its reduction modulo $p^n$ is $D/p^nD$: this
follows by applying the inverse limit to the exact multiplication by
$p$ sequences for $D/p^mD$, whose transition maps are surjective.
Consequently $A_\infty$ is $p$-torsion-free and
$A_\infty/pA_\infty=B$.  Since $B$ is perfect, the strict-$p$-ring
characterization of Witt vectors gives the canonical isomorphism
\[
 A_\infty\simeq W(B).
 \tag{L2.3}
\]
For this characterization see Scholze, Proposition~5.13(i) and
Remark~5.14, or Bhatt, \emph{Lecture II: $\delta$-rings},
Proposition~3.2.  All these identifications commute with restriction
to smaller formal opens.  Indeed, a principal formal open defined by
$f$ has reduction $B_f$ after perfection, and its completed ring is
the unique strict $p$-ring lifting $B_f$; thus the transition maps are
the Witt lifts of the maps on the perfected special fibre.

Define
\[
 D_B=\bigl(W(B)\otimes_{\mathbb Z_p}\mathcal O_C\bigr)^{\wedge}_p,
 \qquad R_B=D_B[1/p].
\]
The $\mathcal O_C$-module $D_B$ is torsion-free: its uncompleted
version is flat over $\mathcal O_C$, $p$-adic completion preserves
$p$-torsion-freeness, and every nonzero scalar of $\mathcal O_C$
divides a power of $p$.  Hence $D_B$ is flat over the valuation ring
$\mathcal O_C$ (also see Scholze, Lemma~5.3(i)).  Its reduction modulo
$p$ is $B\otimes_{\mathbb F_p}\mathcal O_C/p$, and Frobenius induces
an isomorphism $D_B/p^{1/p}\longrightarrow D_B/p$: on the first
factor it is the automorphism of $B$, and on the second it is the
perfectoid Frobenius isomorphism for $\mathcal O_C$.  Thus $D_B^a$
satisfies Definition~5.1(ii) of Scholze, and $R_B$ is perfectoid by
Theorem~5.2.  Under tilting it corresponds to
\[
 R_B^\flat=
 \bigl(B\otimes_{\mathbb F_p}\mathcal O_{C^\flat}\bigr)^{\wedge}[1/\varpi^\flat],
 \tag{L2.4}
\]
where the completion is with respect to any pseudouniformizer.  To
check (L2.4), reduce both integral almost rings modulo matching
pseudouniformizers and use
$\mathcal O_C/p\simeq\mathcal O_{C^\flat}/p^\flat$; their reductions
are the same $B$-algebra.  The uniqueness part of Theorem~5.2 gives
the identification, compatibly with maps in $B$ and with valued-field
automorphisms of $C$.  No choice of a coordinate system or of a
semisimple Frobenius action is involved.

For precision about integral structures, throughout this construction
$\operatorname{Spa}(R_B,R_B^+)$ uses for $R_B^+$ the integral closure
of $D_B$ in $R_B$.  On the characteristic-$p$ side use the integral
closure of $(B\otimes\mathcal O_{C^\flat})^{\wedge}$ in $R_B^\flat$.
These are corresponding plus rings under Scholze's Lemma~6.2: open
integrally closed subrings are determined by their images modulo
the scalar ideal $\mathfrak m_C R_B^\circ$ (respectively
$\mathfrak m_{C^\flat}(R_B^\flat)^\circ$); the two images here
are the integral closures of the same image of
$B\otimes_{\mathbb F_p}\overline{\mathbb F}_p$.  Taking an integral
closure imposes no new condition on valuations already bounded by
the original ring.  Thus (L2.4) identifies the required affinoid
pairs, not just their rings after inverting a pseudouniformizer.

The affinoid perfectoid spaces attached to the $R_B$ glue to a
perfectoid space $\mathcal X_\infty$.  For an explicit check on the
open domains, a principal formal open $D(f)$ has analytic domain
$|f|=1$.  In $W(B)$, the image of $f$ differs from its Teichm\"uller
lift $[\bar f]$ by an element of $pW(B)$, so the same domain is
$|[\bar f]|=1$.  Its rational localization is the completed base
change of $W(B_{\bar f})$.  Finite principal refinements of overlaps
therefore give precisely the same open affinoid domains on both
sides of tilting.  More explicitly, take a finite
formal affine cover of $\mathfrak X$.  The associated affinoid domains
cover $(\mathcal X_C)^{\mathrm{ad}}$ by properness, and their inverse
images under every $\varphi$ are themselves.  On each such domain
$R_B$ is the completion of the direct limit of the affinoid coordinate
rings in the $\varphi$-tower.  A compatible family of bounded
continuous valuations on those rings gives a valuation bounded on
the common integral direct-limit ring.  It is therefore continuous
for the $p$-adic topology and extends uniquely to its completion.  This identifies
the underlying space with the inverse limit of the affinoid spaces;
the induced residue-field maps have dense image.  The same argument
on overlaps glues the identifications.  Hence
\[
 \mathcal X_\infty\ \sim\
 \varprojlim_{\varphi}(\mathcal X_C)^{\mathrm{ad}}
 \tag{L2.5}
\]
in the sense of Scholze's Definition~7.14.

By (L2.4) its tilt is the perfection of the generic fibre of the
constant formal model
$\mathcal X_{\mathbb F_p}\times_{\mathbb F_p}\mathcal O_{C^\flat}$.
For this identification, apply Huber's result first to the
proper noetherian model
$\mathcal X_{\mathbb F_p}\times_{\mathbb F_p}\mathbb F_p[[u]]$,
where $u$ maps to a pseudouniformizer of $C^\flat$, and then base change
to $C^\flat$.  Formal generic fibres commute with that base change on
the displayed affine rings.  This avoids applying Huber's noetherian
statement directly over $\mathcal O_{C^\flat}$.  The generic fibre is
therefore
$(\mathcal X_{\mathbb F_p}\times_{\mathbb F_p}C^\flat)^{\mathrm{ad}}$.
Perfection induces an equivalence of \'etale topoi: use universal
homeomorphism invariance, or Scholze's Corollary~7.19 on the relative
Frobenius tower.  Its finite-coefficient cohomology is therefore that
of $\mathcal X_{\overline{\mathbb F}_p}$, by algebraic--adic comparison
and proper base change (SGA~4, XII, Theorem~5.1 and Corollary~5.2).
Combining this with tilting of \'etale sites and Corollary~7.18 proves
(L2.1).  The construction, reduction, and tilting are functorial for
the semilinear action of $G_{\mathbb Q_p}$ on $C$ and $C^\flat$.
On the constant characteristic-$p$ model the resulting action on
cohomology is exactly the action on the residue field, by proper base
change.  This proves the asserted Galois equivariance.

It remains to justify the rational direct-summand assertion; it is
not an interchange of an inverse limit with a filtered colimit.
Write
$M_i=H^i(\mathcal X_{\overline{\mathbb Q}_p},\mathbb Z_\ell)$,
a finitely generated $\mathbb Z_\ell$-module, and $T=\varphi^*$.
For any finitely generated $\mathbb Z_\ell$-module $M$ with
endomorphism $T$, the sequence $T^{n!}$ converges in
$\operatorname{End}_{\mathbb Z_\ell}(M)$ to an idempotent $e$.
Indeed, on each finite module $M/\ell^aM$ the images of powers of $T$
stabilize, $T$ is a permutation of the stable image, and a sufficiently
large factorial kills the transient part and is divisible by the
order of that permutation.  The resulting idempotents are compatible
as $a$ varies.  Thus
\[
 M=M^{\mathrm{u}}\oplus M^{\mathrm{n}},\quad
 M^{\mathrm{u}}=eM,\quad
 \mathop{\rm colim}_T(M/\ell^aM)=M^{\mathrm{u}}/\ell^aM^{\mathrm{u}}.
 \tag{L2.6}
\]
Every operator commuting with $T$, including the Galois action,
commutes with $e$.  On $M^{\mathrm{u}}$ the map $T$ is invertible,
and on the other summand it is topologically nilpotent.

The universal-coefficient exact sequence is
\[
 0\longrightarrow M_i/\ell^aM_i\longrightarrow
 H^i(\mathcal X_{\overline{\mathbb Q}_p},\mathbb Z/\ell^a)
 \longrightarrow M_{i+1}[\ell^a]\longrightarrow0.
\]
Take the exact filtered colimit under $T$ and use (L2.1) and (L2.6).
The resulting rightmost groups are killed by a fixed power of $\ell$,
independent of $a$, since the torsion subgroup of $M_{i+1}$ is finite.
The leftmost groups are $M_i^{\mathrm{u}}/\ell^aM_i^{\mathrm{u}}$,
with surjective transition maps.  Taking the inverse limit and then
inverting $\ell$ therefore yields
\[
 H^i(\mathcal X_{\overline{\mathbb F}_p},\mathbb Q_\ell)
 \simeq M_i^{\mathrm{u}}\otimes\mathbb Q_\ell
 \subseteq M_i\otimes\mathbb Q_\ell.
 \tag{L2.7}
\]
The inclusion is a $G_{\mathbb Q_p}$-stable direct summand by the
idempotent $e$.  This finishes the proof.
\end{proof}

\begin{theorem}[NX-L03: bad-prime splitting for proper flat $\Lambda$-schemes]
\label{NX-L03}
Let $X$ be proper and flat of finite type over $\mathbb Z$, with a
$\Lambda$-structure.  For every prime $p$ and every $\ell\ne p$,
\[
 H^i(X_{\overline{\mathbb F}_p},\mathbb Q_\ell)
 \simeq\bigoplus_j V_{p,i,j}\otimes\mathbb Q_\ell(-j),
 \tag{L3.1}
\]
where each $V_{p,i,j}$ has finite image under
$G_{\mathbb F_p}$ and is semisimple.  In particular geometric
Frobenius is semisimple and its eigenvalues have the form
$\xi p^j$, with $\xi$ a root of unity and $j\in\mathbb Z$.
After a finite extension of $\mathbb F_p$, every cohomology group is
a direct sum of integral Tate twists.  The bounded derived
$\mathbb Q_\ell$-representation $R\Gamma_c(X_{\overline{\mathbb F}_p},\mathbb Q_\ell)$
likewise splits, noncanonically, into those summands in their respective
degrees.  Consequently every factor of $R_{X,S}$ in ZETA\_TWO~ZT04
is cyclotomic integral-Tate for this proper class, including at bad
primes.
\end{theorem}

\begin{proof}
Apply NX-L02 to $X\times\mathbb Z_p$ and $\varphi=\psi_p$.
Borger's Theorem~6.1 and Corollary~6.2 apply to the proper, hence
separated, scheme $X$: its generic-fibre cohomology is potentially
a direct sum of integral Tate twists, and its global Galois action
is abelian.  Over $\overline{\mathbb Q}_\ell$ it is thus a direct
sum of characters $\epsilon\chi_\ell^{-j}$ with $\epsilon$ of
finite order.  One can see semisimplicity without enlarging the
coefficient field by splitting after the finite-index subgroup in
Corollary~6.2 and averaging an equivariant projection over the finite
quotient.  The $G_{\mathbb Q_p}$-stable direct summand supplied by
NX-L02 is semisimple of the same form.  Its inertia action is trivial
because it is the cohomology of a variety over $\mathbb F_p$.
Therefore only unramified constituents occur, and geometric
Frobenius on each is $\epsilon(F_p)p^j$.  This is (L3.1), with the
finite characters grouped over $\mathbb Q_\ell$.

For the assertion about the derived object, use the elementary
hereditary property of finite-dimensional representations of an
invertible operator over $\mathbb Q_\ell$.  They are torsion modules
over the principal ideal domain $\mathbb Q_\ell[T,T^{-1}]$;
$\operatorname{Ext}^r=0$ for $r\geq2$.  In the continuous subcategory
relevant here the eigenvalues are $\ell$-adic units, and the same
finite-length module calculation applies.  The truncation triangles
of a bounded complex consequently split successively.  This gives
a noncanonical direct sum of its cohomology objects shifted to their
degrees.  Properness identifies compactly supported and ordinary
cohomology.  Notice that this argument does not invoke a motivic
splitting or the Tate conjecture.

Finally the trace formula of SGA~5, XV, \S3, no.~2, unnumbered theorem,
equation~(1), expresses the local zeta function as
\[
 Z(X_{\mathbb F_p},t)=
 \prod_i\det(1-tF_p\mid H^i(X_{\overline{\mathbb F}_p},\mathbb Q_\ell))^{(-1)^{i+1}}.
\]
Every determinant splits into the claimed factors by (L3.1).  The
canonical local factors of the abelian Artin--Tate generic
representations have the same form, since inertia invariants of a
finite-image representation again have finite-image Frobenius.
Taking their ratio is precisely ZT04.  This proves the final claim.
\end{proof}

\begin{proposition}[NX-L04: proper splitting does not preserve multiplicities]
\label{NX-L04}
Even for a finite flat proper $\Lambda$-scheme, the bad special-fibre
cohomology need not equal the generic-fibre cohomology or its inertia
invariants.  In particular NX-L03 does not remove the finite
correction in ZT04.
\end{proposition}

\begin{proof}
Take $X=\mu_p=\operatorname{Spec}\mathbb Z[t]/(t^p-1)$, with
$\psi_q(t)=t^q$.  The quotient is free of rank $p$ over $\mathbb Z$;
the endomorphisms commute, preserve the relation, and lift Frobenius.
The generic geometric fibre has $p$ points, whereas the reduced
special fibre at $p$ is a single rational point.  Thus generic $H^0$
has dimension $p$ and special $H^0$ has dimension $1$.
Moreover $X_{\mathbb Q}$ has one rational point and the cyclotomic
factor $\operatorname{Spec}\mathbb Q(\zeta_p)$ (for $p=2$ the latter
is a second rational point).  Its generic canonical local factor at
$p$ is $(1-t)^{-2}$ because each number-field component has one
residue-degree-one prime over $p$.  Its actual special-fibre local
factor is $(1-t)^{-1}$.  The ZT04 correction is therefore $1-t$.
This also shows directly that special $H^0$ is not the full
inertia-invariant part of generic $H^0$.
\end{proof}

\begin{corollary}[NX-L05: the vertical elliptic blowup has no Frobenius lift]
\label{NX-L05}
Let $E\subset\mathbb P^2_{\mathbb F_p}$ be a smooth plane cubic and set
$\mathcal B=\operatorname{Bl}_{E}(\mathbb P^2_{\mathbb Z_p})$.
This is a proper flat scheme with generic fibre $\mathbb P^2_{\mathbb Q_p}$,
but it has no $\mathbb Z_p$-endomorphism lifting absolute Frobenius on
its special fibre.  In particular a vertical elliptic blowup of a
Tate generic fibre cannot provide a proper counterexample to NX-L03.
More generally, any proper flat model with
$H^i(\mathcal X_{\overline{\mathbb Q}_p},\mathbb Q_\ell)=0$ and
$H^i(\mathcal X_{\overline{\mathbb F}_p},\mathbb Q_\ell)\ne0$ has no
such Frobenius lift.
\end{corollary}

\begin{proof}
The general assertion is immediate from the direct-summand conclusion
of NX-L02.  We compute the indicated nonzero group for the blowup,
so that the obstruction is explicit.  Locally the centre is cut out
by the regular sequence $(p,F)$, with $F\bmod p$ an equation for $E$.
The blowup charts have equations $F=pu$ and $p=Fv$.  The total blowup
is integral, since its Rees algebra embeds in the polynomial ring
of the integral ambient scheme, and is therefore flat over the DVR
$\mathbb Z_p$.  It is projective, and after inverting $p$ the ideal
is the unit ideal, giving the stated generic fibre.

Its reduced special fibre, which here is its whole special fibre,
is the union of the strict transform $D\simeq\mathbb P^2_{\mathbb F_p}$
and the exceptional divisor $P=\mathbb P(N_{E/\mathbb P^2_{\mathbb Z_p}})$,
a $\mathbb P^1$-bundle over $E$; their intersection is the copy of $E$
which is the section corresponding to the rank-one normal bundle
inside the special fibre.  These assertions also follow immediately
from the two displayed blowup charts; the strict transform is unchanged
because blowing up the Cartier divisor $E$ in $\mathbb P^2_{\mathbb F_p}$
is the identity.

The constant-sheaf sequence for the union $D\cup_E P$ is exact on
all geometric stalks.  Its long cohomology sequence contains
\[
 H^2(D)\oplus H^2(P)\longrightarrow H^2(E)
 \longrightarrow H^3(\mathcal B_{\overline{\mathbb F}_p})
 \longrightarrow H^3(D)\oplus H^3(P)\longrightarrow H^3(E).
\]
All coefficients here are $\mathbb Q_\ell$.  The first arrow is
surjective: the pullback of $H^2(E)$ to $P$ restricts identically along
the section $E\hookrightarrow P$.  Also $H^3(D)=H^3(E)=0$.
For the projective-line bundle $P\to E$, proper and smooth base
change give relative cohomology $\mathbb Q_\ell$ in degree zero,
$\mathbb Q_\ell(-1)$ in degree two, and zero otherwise.  The relative
hyperplane class supplies the latter identification.  In its Leray
spectral sequence the only term of total degree three is
$H^1(E,\mathbb Q_\ell)(-1)$, and no differential can reach or leave
it because $E$ has cohomological dimension two.  Consequently
\[
 H^3(\mathcal B_{\overline{\mathbb F}_p},\mathbb Q_\ell)
 \simeq H^1(E_{\overline{\mathbb F}_p},\mathbb Q_\ell)(-1),
\]
a two-dimensional group.  On the generic fibre
$H^3(\mathbb P^2_{\overline{\mathbb Q}_p},\mathbb Q_\ell)=0$.
NX-L02 therefore excludes a Frobenius lift.  In weights language this
is a weight-three elliptic contribution in the special fibre.  The
base-change inputs used in this computation have the exact locators
SGA~4, XII, Theorem~5.1 and Corollary~5.2, and XVI, Theorem~2.1.
\end{proof}

\begin{corollary}[NX-L06: the corrected proper $\Lambda$-RH conclusion]
\label{NX-L06}
Using the stipulated Alp\"oge theorem A, let $X$ be a proper flat
$\Lambda$-scheme of finite type over $\mathbb Z$.  Then $\zeta_X(s)$
has meromorphic continuation, and every non-real zero or pole has
real part in $\tfrac12\mathbb Z$.  More precisely, its global
Dirichlet--Tate factors have their non-real divisors on lines
$\Re(s)=j+\tfrac12$ with $j\in\mathbb Z$, while its actual local
factors and the finite correction factors in ZT04 have their divisors
on lines $\Re(s)=j$ with $j\in\mathbb Z$.
\end{corollary}

\begin{proof}
Borger, Theorem~6.1 and Corollary~6.2, give abelian Artin--Tate
cohomology of $X_{\mathbb Q}$, so on a good open its virtual zeta
factor is a finite product of integral shifts of primitive Dirichlet
$L$-functions with integer exponents.  Restoring the omitted primes
is exactly ZETA\_TWO~ZT04.  NX-L03 makes each actual and correction
factor a product of $(1-\xi p^{j-s})^{\pm1}$, whose divisor has
$\Re(s)=j$.  The supplied Alp\"oge \emph{Remark}, step~20, p.~4,
and the endpoint declarations
\texttt{LFunction\_eq\_zero\_re} and
\texttt{LFunction\_ne\_zero\_of\_half\_lt\_re} in
\texttt{ALPOGE\_Main.lean}, supply the Dirichlet zero-location and
nonvanishing statements.  Together with the classical functional
equations and their real trivial divisors, they place the non-real
divisors of the global shifted factors on $j+\tfrac12$, as in
ZETA\_TWO~ZT01.  Cancellation in a finite meromorphic product removes
divisor support and cannot create new support.  This proves the claim.
No Masullo input enters this corollary.
\end{proof}

\paragraph{Exact scope of the resolution.}
NX-L01 refutes the conjecture literally asked, since separatedness
was not one of its hypotheses.  NX-L02--NX-L03 answer the proper flat
question positively on actual cohomology, including semisimplicity;
NX-L04 shows why equality with the whole generic cohomology or with
its inertia invariants would be a different and false assertion.
A separately proposed \emph{separated, nonproper} strengthening is
not settled by these arguments: the counterexample is nonseparated,
and the formal generic-fibre comparison used in NX-L02 relies on
properness.  No $\Lambda$-equivariant compactification is asserted.
This limitation does not leave the literal conjecture undecided.

\begin{proposition}[NX-L07: cyclotomic local factors do not imply a $\Lambda$-structure]
\label{NX-L07}
The integer ring $\mathbb Z[\zeta_5]$ has cyclotomic integral-Tate
local zeta factors but admits no Frobenius lift at $5$, hence no
$\Lambda$-structure.  Thus the prompt's local-factor observation for
abelian integer rings is valid without asserting that all those
rings are $\Lambda$-rings.
\end{proposition}

\begin{proof}
Every nonzero ideal of the ring of integers $\mathcal O_K$ of a
number field meets $\mathbb Z$ nontrivially: if $0\ne x$ lies in the
ideal, its monic minimal polynomial gives a nonzero integer constant
term in the ideal $(x)$.  Therefore every unital $\mathbb Z$-algebra
endomorphism of $\mathcal O_K$ is injective.  It extends to an
endomorphism of $K$, hence to an automorphism since $K$ is
finite-dimensional over $\mathbb Q$, and preserves its integral
closure $\mathcal O_K$ in both directions.  Its reduction modulo
any integer is consequently an automorphism.

Apply this to $K=\mathbb Q(\zeta_5)$.  Reducing the cyclotomic
polynomial modulo $5$ gives
$\Phi_5(T)=(T-1)^4$ in $\mathbb F_5[T]$, so
\[
 \mathbb Z[\zeta_5]/5\mathbb Z[\zeta_5]
 \simeq\mathbb F_5[\varepsilon]/(\varepsilon^4).
\]
Its Frobenius kills the nonzero element $\varepsilon$, and is not an
automorphism.  It cannot be the reduction of a ring endomorphism.
On the other hand, each local zeta factor of $\operatorname{Spec}\mathcal O_K$
is a product of $(1-t^f)^{-1}$ over the residue degrees $f$ of the
primes above the rational prime.  Each such factor splits into
$(1-\xi t)^{-1}$ over the $f$th roots of unity.  These have integral
Tate index $j=0$, proving the local-factor claim independently of
the nonexistent $\Lambda$-structure.
\end{proof}

\subsection{Placement in Grothendieck's frameworks}

NX-L01 is a statement about flat nonseparated schemes over
$\operatorname{Spec}\mathbb Z$, obtained by Zariski gluing.  Its
local zeta factor is a multiplicative invariant of point counts;
the weight-one factor is the realization of $h^1(E)$ and cannot be
replaced by integral Tate weights.  The full half-integer-line
statement is not refuted: those two inverse roots have weight one,
so their local divisors lie on $\Re(s)=1/2$.

NX-L02 and NX-L03 concern actual $\ell$-adic realizations, including
semisimplicity, and hence strengthen a mere assertion in a
Grothendieck group.  Their local determinants are computed in the
$\lambda$-ring $K_0$ of finite-dimensional representations, using
exterior powers.  This is the formalism of SGA~6, V, \S2 and
\S7, (7.1.1), used already in ZETA\_TWO~ZT08 and SV01.  The sheaf
operations and proper base change live in SGA~4, XII, Theorem~5.1
and Corollary~5.2; compact supports are treated in XVII, \S\S5.1--5.3;
the zeta determinant is SGA~5, XV, \S3, no.~2, equation~(1).
The perfectoid step is a later comparison theorem, not a theorem
silently ascribed to EGA, SGA, or FGA.

For the archive connection, the supplied predecessor places exterior
and symmetric powers and characteristic determinants at cote~60,
scans~3 and~5 (ZETA\_TWO, \emph{IMPLICATIONS}, \S5, placement table);
its scan~7, archival p.~6, records a reciprocal characteristic-polynomial
identity, while scan~9, archival p.~8, records Weil and dual/Tate
normalizations.  These exact locators are carried from the checked
predecessor, not newly inspected scans in this run.  DP\_MULT~DP11
and its intake identify the binomial/identity-Adams special case;
that case does not encompass NX-L01, whose Frobenius lifts are
nonidentity.  DP\_MULT's reported cote~153, PDF p.~3/archive p.~2,
concerns binomial polynomials, not a bad-fibre cohomology theorem.
No motive-level or cycle-level decomposition is asserted by (L3.1).
